\subsection{Spaced Repetition Implementation}
\label{sec:impl_sr}

The spaced repetition subsystem translates the SM-2 algorithm and its extensions (described in Section~\ref{sec:system_sm2}) into a production-grade scheduling service. The implementation lives in \path{abridgeai/features/spaced_repetition/}. Three layers carry the scheduling cycle itself --- the pure SM-2 core, the review service that composes it into a transaction, and the background worker that raises reminders from the resulting due dates --- and two read-side surfaces are built on the state they produce: the lesson-unlock gate and the instructor analytics queries. Each is described in turn below.

\paragraph{SM-2 core}
The algorithmic core is implemented as a set of pure functions in \path{features/spaced_repetition/sm2/}:

\begin{itemize}
    \item \path{ef_update.update_ef(ef_old, q, n, alpha=0.6, *, positive_delta_scale=1.0)} — computes the new Easiness Factor using the standard SM-2 delta formula with calibration dampening. During the first three repetitions ($n \leq 3$), positive EF deltas are multiplied by $\alpha = 0.6$ to slow early growth and mitigate over-confidence from limited evidence. Negative deltas are not dampened so the forgetting signal is preserved. The result is clamped to $[EF_{\min}, EF_{\max}] = [1.3, 2.5]$ and rounded to four decimals; the second multiplier \texttt{positive\_delta\_scale} carries the guess-channel dampening described below.
    \item \path{scheduler.next_interval_days(*, ef, n, q, prev_interval)} — computes the next review interval in days following the SM-2 recurrence: 1 day for the first review, 6 days for the second, and $\max(1, \text{round}(\text{prev\_interval} \times \text{ef}))$ for subsequent reviews. A grade of $q < 3$ resets the interval to 1 day unconditionally.
    \item \path{scheduler.apply_interval_ceiling(interval_days, *, max_interval_days, retire_beyond)} — bounds the interval at the configured ceiling and reports whether that bound was exceeded, returning \texttt{(interval, retired)}. Applied after jitter so the ceiling cannot be overshot by the jitter itself. Retirement is reported only when the interval strictly exceeds the bound, so a ceiling set to a desired resting cadence yields that cadence rather than a retirement.
    \item \path{scheduler.apply_jitter(interval_days, fraction=0.1)} — adds $\pm 10\%$ random jitter to the computed interval to prevent review pile-up when a cohort of students all enroll on the same day. Jitter uses float arithmetic (\texttt{round(interval * (1 + epsilon))}) to avoid the dead zone that integer truncation produces for short intervals.
    \item \texttt{q\_derivation} — derives the $Q$ value from objective response signals (correctness, hint usage, and time ratio $\rho = T_{\text{actual}} / T_{\text{exp}}$) as specified in Section~\ref{sec:q_derivation}.
\end{itemize}

The four modules above are pure functions with no database or network dependencies, making them straightforwardly unit-testable. The test suite in \path{tests/unit/} covers boundary conditions including $\rho$ at 0.5, 1.0, and 2.0, $n = 0$ through $n = 4$, and $q \in \{0, 1, 2, 3, 4, 5\}$. The \path{sm2/} package additionally holds \path{lesson_unlock.py}, which is \emph{not} pure --- it reads lesson configuration and per-card state from the database --- and is therefore exercised by the integration suite (\path{tests/integration/test_lesson_unlock.py}) rather than the pure-function unit tests above. It is described under \textit{Lesson unlock} below.

\paragraph{Review service}
The review service (\path{features/spaced_repetition/services/review.py}) orchestrates a card review transaction:

\begin{enumerate}
    \item Load the question's scheduling metadata ($T_{\text{exp}}$, parent quiz, teacher-configured initial EF) through the quizzes public API, then load — or lazily initialise — the \texttt{StudentCardState} row keyed by the composite primary key (\texttt{student\_id}, \texttt{question\_id}).
    \item Derive $Q$ from the submitted response signals.
    \item Call \texttt{update\_ef} and \texttt{next\_interval\_days} to compute the new EF and interval.
    \item Apply jitter via \texttt{apply\_jitter} and compute the absolute \texttt{due\_at} timestamp. On the failure path ($q < 3$) the jittered interval is bypassed: the card resets to $n = 0$, interval 1, and \texttt{due\_at} is pushed out by a configurable cooldown (default 24 hours) that the quiz router reads back as \texttt{retry\_available\_at} to enforce a per-card retry block.
    \item Apply the interval ceiling. Where retirement is enabled and the bound is exceeded, \texttt{due\_at} is set to \texttt{NULL} rather than to a timestamp; every due-card read already filters on \texttt{due\_at IS NOT NULL}, so a retired card leaves the queue without any read requiring knowledge of a new state, and without its history being altered.
    \item Mutate the \texttt{StudentCardState} row in place and insert an immutable \texttt{CardReview} row — the full before/after audit record (\texttt{ef\_before}/\texttt{ef\_after}, \texttt{interval\_before}/\texttt{interval\_after}, \texttt{n\_before}/\texttt{n\_after}, $\rho$, and the derived $Q$) that the analytics queries read.
\end{enumerate}

The service takes no row-level lock. Both writes are issued inside the caller's transaction and the caller owns the \texttt{commit()}, so a failure at any step rolls the card state and the audit row back together. Side effects are deliberately deferred past that boundary: on $q = 0$ the service does not dispatch a remediation notification at flush time but returns a \texttt{CardFailedEvent} in the result's \texttt{pending\_events} list, which the caller dispatches only after the commit succeeds. This closes a ghost-notification race in which a rolled-back review still fired a learner-visible side effect.

\paragraph{Assessment quizzes open no cards}
Card creation is driven by the quiz answer flow, which calls the review service on a student's first answer to each question. A quiz carries a \texttt{feeds\_spaced\_repetition} flag, defaulting to true, that governs whether that call may \textit{initialise} state for a question the student has not met before. The decision is taken in the quizzes feature rather than inside the review service: the caller consults the spaced-repetition public API for existing card state and suppresses the call only when the quiz is an assessment \textit{and} no card exists, leaving the review service's contract unchanged for every other caller.

The asymmetry is deliberate. Suppressing the creation of new cards prevents an examination from populating a student's review queue with every question it asked; suppressing the grading of existing cards as well would allow an assessment to shield an item from the evidence of having been failed, which inverts the purpose of assessing it. Like the other SM-2 parameters on a quiz, the flag is excluded from the set of fields editable after publication, so it cannot change beneath a live or completed attempt.

\paragraph{Lesson unlock}
Unlock is implemented in \path{sm2/lesson_unlock.py} as the combined gate specified in Section~\ref{sec:individualised_modeling}, evaluated over three independent conditions --- shown together with the enrolment checks that precede them and the content boundary they protect in Figure~\ref{fig:lesson_unlock_gate} of Appendix~\ref{appendix:algorithm}:

\begin{itemize}
    \item \textbf{Prerequisite gate} — every lesson listed in \texttt{lesson\_prerequisites} must itself be eligible. The check recurses through the prerequisite graph and is cycle-safe: an $A \rightarrow B \rightarrow A$ edge short-circuits with a warning log and is treated as satisfied, so a malformed graph never bricks a learner out of their own course.
    \item \textbf{EF gate} — the aggregation query (\path{queries/sql/lesson_unlock.sql}) walks \texttt{lessons} $\rightarrow$ \texttt{modules} $\rightarrow$ \texttt{module\_items} $\rightarrow$ \texttt{quizzes} $\rightarrow$ \texttt{quiz\_questions} and returns three primitives in one shot: the count of cards whose stored EF is at least the lesson's \texttt{ef\_min\_unlock}, the total card count, and a capped JSONB array describing the cards that fall below the floor. The gate passes when $\text{passing} / \text{total} \geq \tau$, where $\tau$ is the lesson's \texttt{tau\_unlock}. A \texttt{LEFT JOIN} onto \texttt{student\_card\_state} means a card the student has never reviewed counts as $EF = 0$ and therefore blocks. A lesson with zero cards bypasses this gate rather than dividing by zero.
    \item \textbf{Interview gate} — when the lesson sets \texttt{requires\_interview\_pass}, the student must additionally hold a completed \texttt{interview\_sessions} row with \texttt{pass\_verdict = TRUE} against the interview config attached to the lesson's parent module.
\end{itemize}

Two points of terminology are worth stating precisely, because the quantity this gate computes is not the one its name might suggest. The ratio above is the EF-\emph{coverage} proportion --- the fraction of cards clearing a floor --- which Section~\ref{sec:performance_metrics} names \textit{Progression Readiness}. It is not the Knowledge Retention Estimate $\hat{R}_{s,\ell}$, which is a different quantity (the mean normalised EF, $\frac{1}{|\mathcal{C}_{\ell}|}\sum (EF - 1.3)/(2.5 - 1.3)$), computed by a different query (\path{queries/sql/kr_estimate.sql}), over a differently-derived card set (quizzes linked through \texttt{quiz\_source\_lessons} rather than through \texttt{module\_items}). The two are complementary, not interchangeable: a student may sit well above the mean while a single card below the floor holds the gate shut at $\tau = 1.0$. Nor does the gate consult the scheduling interval at any point --- there is no interval threshold in the condition, only the EF floor and the coverage proportion.

There is likewise no unlock service and no unlock event. The gate is a pure read-side computation: \texttt{check\_lesson\_unlock} is invoked by the learner router on each request that needs the verdict, and nothing is emitted or persisted when a threshold is crossed. Because the traversal is recursive and per-lesson, the result is memoised behind the \texttt{LESSON\_UNLOCK} cache key with a 60-second TTL, and the cache invalidator additionally drops that key whenever \texttt{student\_card\_state}, \texttt{card\_reviews}, or \texttt{lessons} is written --- so a review that crosses the threshold is reflected on the learner's next read rather than after the TTL elapses. Alongside the binary verdict the function returns the supporting evidence the UI renders: the current and required ratios, the passing and total card counts, the blocking cards with their current EF and source-chunk references for targeted remediation, and a human-readable next-step estimate (\texttt{``Review 3 more cards to reach the unlock threshold.''}).

\paragraph{Background worker}
An \texttt{arq} cron worker (\path{features/spaced_repetition/workers/scan_due_cards.py}) runs hourly --- registered in \path{abridgeai/workers/arq_app.py} as \texttt{cron(scan\_due\_cards\_task, minute=0)} --- finding students whose cards are due, meaning \texttt{due\_at} has passed and the existing backlog is included, and dispatching one in-app notification per student through the notifications subsystem. Per-lesson rows are folded into a single notification that names the lessons involved, so a student with cards across three lessons receives one ping rather than three.

The due predicate is deliberately backward-looking. An earlier implementation used a forward window (\texttt{due\_at BETWEEN now AND now + 1h}), which counted only cards about to \emph{become} due and silently excluded everything already overdue --- so a student with a genuine backlog, the one case the reminder exists to serve, received no reminder at all, and the figure in any reminder that did fire disagreed with every read surface in the product. The scan now mirrors the cards-due page exactly: the same joins, the same \texttt{due\_at <= NOW()} predicate, the same reviewability and soft-delete filters. The count in the notification is therefore the count on the page, by construction.

The hourly cadence keeps newly-due cards surfacing promptly, but on its own it would re-notify a standing backlog once an hour, pinging a student who has fallen behind twenty-four times a day. Two guards prevent that. A per-student cooldown skips anyone who already received a spaced-repetition notification inside the configurable window (default 24 hours), so the scan stays responsive to \emph{new} backlog without re-announcing old backlog. A trigger gate then restricts pings to students holding at least one due card from a quiz with \texttt{reminders\_enabled}; a student whose entire backlog comes from opted-out quizzes is never pinged. Crucially the flag gates only \emph{who} is notified; it is not applied as a count filter, so a pinged student is told their full backlog rather than the opted-in slice of it --- preserving the count contract above. Dispatch is in-app only: the email channel is deliberately not wired for due-card reminders.

\paragraph{Robustness extensions beyond textbook SM-2}
Four defences were added on top of the algorithm as specified, each closing a way the textbook formulation misbehaves against real learner input rather than idealised grades.

\textit{Guess-channel dampening.} A correct answer on a closed-option format carries a probability of having been a lucky guess that free recall does not. Before the EF update, the service reads a per-format guess probability $P(G)$ --- $0.5$ for true/false, $1/N$ for an $N$-option multiple-choice item, and $0$ for short-answer, fill-in-the-blank, numerical, code, matching and ordering items --- and passes $1 - P(G)$ as the positive-delta scale. A guessable format therefore grows EF more slowly than genuine retrieval, so a fast lucky click cannot balloon the interval. Crucially the mechanism does not touch $Q$: the thesis Q-derivation is unchanged and the derived value is still recorded verbatim in \texttt{card\_reviews}; only the rate of EF growth is softened. Like the calibration weight $\alpha$, the scale applies to positive deltas only, leaving the forgetting signal at full magnitude. Free-recall formats pass $1.0$ and are unaffected.

\textit{EF ceiling.} Canonical SM-2 imposes no upper bound on the Easiness Factor, but every EF consumer in this system assumes one: the initialisation value is $2.5$, the database CHECK constraint admits $[1.3, 2.5]$, and the KR normalisation divides by $EF_{\max} - EF_{\min} = 2.5 - 1.3$. Left uncapped, a run of perfect reviews drifts EF past $2.5$ and the retention estimate silently exceeds $1.0$ --- a retention figure above 100\%. Clamping the update at $EF_{\max} = 2.5$ keeps $EF \in [1.3, 2.5]$ and hence $\hat{R}_{s,\ell} \in [0, 1]$ by construction rather than by convention.

\textit{Response-time clamp.} $T_{\text{actual}}$ is client-reported and therefore untrusted. Because \texttt{derive\_q} buckets $\rho$ and floors at $Q = 3$ for a correct unhinted answer, $\rho = 1.01$ and $\rho = 60$ are already indistinguishable to the model --- so the service clamps $T_{\text{actual}}$ to three times $T_{\text{exp}}$ at zero cost in fidelity, keeping a student who left the tab open over lunch, a skewed clock, or a hand-crafted payload out of the stored timings that analytics later average. The client cannot apply this ceiling itself: the student-facing question payload deliberately withholds \texttt{expected\_response\_time\_ms}, since exposing it would leak how long the teacher expects the question to take. Negative values floor at zero as defence in depth behind the schema's own \texttt{ge=0}.

\textit{Initial-EF clamp.} A teacher may set a per-quiz \texttt{initial\_ef} that the quiz settings accept without range validation. Seeding a new card from an out-of-range value would violate the \texttt{ck\_student\_card\_state\_ef\_range} constraint and surface as an \texttt{IntegrityError} on the student's first review --- so the value is clamped into $[1.3, 2.5]$ at card creation, where the failure is invisible to the learner rather than fatal to their attempt.

\paragraph{Analytics dashboards}
The \path{features/spaced_repetition/routers/teacher.py} router exposes instructor-facing analytics endpoints backed by SQL queries in \path{queries/sql/}:
\begin{itemize}
    \item \texttt{at\_risk.sql} — students whose average EF has dropped below a threshold, indicating systematic forgetting.
    \item \texttt{compliance\_rate.sql} — fraction of due reviews completed on time per student per lesson.
    \item \path{class_kr_distribution.sql} — histogram of KR scores across the class for a given lesson.
\end{itemize}

\paragraph{The cycle end to end}
Taken together the pieces form a closed cycle. A graded answer yields $Q$; $Q$ drives the ease-factor update; the updated ease factor and the repetition count give the next interval; jitter spreads that interval so a cohort enrolled on one day does not return on one day; and the resulting \texttt{due\_at} is what the background scan later reads to raise a reminder. An incorrect answer short-circuits the middle of that chain: a grade below three resets the interval and the repetition count regardless of the ease factor and holds the card behind a cooldown, so a lapse returns the card to the front of the queue rather than merely shortening its next gap.

Unlock sits outside this cycle rather than inside it. The card state written at the end of each pass is the state the unlock aggregation reads, but nothing pushes the verdict forward at write time --- the gate is recomputed from that state whenever a learner request needs it, which is what lets the same query answer both ``may I proceed?'' and ``what exactly is holding me back?'' from one source of truth. Writes to \texttt{student\_card\_state} and \texttt{card\_reviews} only invalidate the memoised verdict; they never compute a new one. The two can therefore never disagree about what the learner has demonstrated, because there is only ever one evaluation and it happens on read.
